% !TeX root = ../papers/nc-p3.tex
\section{階数15の種数3四曲線上の全通常点}
\label{sec:rank-fifteen-octic-all}
\begin{theorem}\label{thm:rank-fifteen-octic-all}
標数零とする。$\xi^2+\xi+1=0$ の両水準で、階数15の両表示の
\emph{全通常点}は標準Koszul左右AS正則である。各始点で道空間は
四つの四次中心元が生成する多項式環上の階数256の自由加群になる。
\end{theorem}
\begin{proof}
まず109で $t$ が八次曲線の通常域に還元される場合を扱う。
$k=\mathbb F_{109}$、$F=t^8+14t^4+1$ とし、$k(t)[Y]/(Y^2-F)$ 上で
$j=8,41$ と両表示を取る。$F$ は109でも平方自由である。
\path{computations/rank_fifteen_octic_modular.py} は分母をモニックにした
$(a+bY)/d$ で消去する。道を長さ5まで延ばし、各射影の係数の極と実際の
階数小行列式の例外を記録する。四次中心の交換子にも同じ判定を行う。
有理関数の非零判定を一変数で行うには $a^2-Fb^2$ と $d$ を記録すればよい。
得られる道次元は $1,4,10,20,35,56$、中心次元は $(3,2,0,0)$ である。

\path{computations/verify_rank_fifteen_octic_ordinary.py} は文字1の中心二本を固定し、
文字0の三本の核基底のうち、第一・第三の対、または係数 $(1,3,5),(2,7,11)$ の対を使う。
始点0,1のそれぞれで商の道に一文字を加え、末尾の二次関係と長さ4での
中心関係を順次課す。分解する有理剰余で独立な行を選び、行列式は上の関数体で取る。
各段の選択小行列式はピボットの積である。前段までの関係・中心式・射影が
正則でこの小行列式が単元なら、選択列での逆をCramerの公式で取れる。
従って各始点で商次元の標準上界
\[
 1,4,10,20,31,40,44,40,31,20,10,4,1,0
\]
が長さ13まで得られる。二つの選択の例外多項式は、分母とノルム分子の
各平方自由部分の最小公倍数でよく、両方の例外に入る因子のみを残す。

始点0,1の残余の和集合を取り、既知の非通常因子と一次因子を除く。
$(\text{表示},j)=(0,8),(0,41),(1,8),(1,41)$ の順に
それぞれ $23,17,28,21$ 個の相異なる既約因子を得る。各因子の根と $Y$ の両シートについて
剰余体上で関係、道、中心、四中心商を再計算する。非分解の場合は
Frobeniusで交換される二点を一組として扱える。全てで全四始点の上界を得る。
一次の通常剰余は命題\ref{prop:rank-fifteen-octic-residues}で認証済みである。
$t\mapsto-1/t$ の窓移動で残る二始点へ、$j\mapsto-j$ の窓移動で残る定数根へ移る。
始点ごとに四式の選択が異なっても、良い四式の条件は各始点で同じ係数空間の
非空開条件であり、代数閉体上で四始点に共通な選択がある。
定理\ref{thm:rank-fifteen-octic-center}の中心次元と道次元が剰余でも保たれるので、
中心が持ち上がり、通常な持上げには全次数の自由性と正則性が従う。

次に $K=\pm1$ に還元される場合を扱う。$D=\mathbb Z_{(109)}[i,\sqrt3]$、
$j=(\sqrt3+i)/2$ とし、局所変数に $u=Y$ を取る。境界点の
$\alpha=(V-1)/(U-1)$ と $U,V$ の全符号は節\ref{sec:rank-fifteen-octic-boundary}の通りである。
$F'(\alpha)$ は $D$ の単元であり、$F(t)=u^2$ は $t(0)=\alpha$ から
$D[[u]]$ に解ける。\path{computations/rank_fifteen_octic_integral.py} は
標数零の八次曲線の関数体で、長さ4,5の経路基底と四次中心の核を正確に求める。
文字0の第一・第三と文字1の二本を選び、各中心元の全四位置で
一つの共通分母を払う。交換子は通常域の定階数の経路空間で一般点において零なので、
これらは分母を払った後も通常域全体で中心となる。

\path{computations/verify_rank_fifteen_octic_trace_disks.py} は
$D[[u]]/(u^{24})$ で二次関係と四中心式を展開する。
$1,i,\sqrt3,i\sqrt3$ の四成分を有理数係数の打切り多項式で保持し、
単元判定は定数項を109上の全四埋込み $(i,\sqrt3)=(\pm33,\pm49)$ で行う。
商の各文字ブロックで単元ピボットを順に取り、必要な階数に達する前に
単元が尽きた場合、残る全係数の定数項が数体上で厳密に零であることを確かめて
$u$ で割る。通常な持上げでは $u\ne0$ なので得られる式は関係として使える。
各文字で割った回数の最大値だけ次段の精度を下げ、各段の自由語は商を張る。
二位置移動が $(U,V)$ の両符号を同時に反転するため、第一の符号を固定して
第二の両符号と全四始点を取ればよい。全32組で長さ13に達しても
精度が正に残り、上記の標準上界を得る。この計算は二次極限で増加した経路を
高次関係の飽和で取り除く。

以上で109における通常剰余と $K=\pm1$ の全円板を尽くした。残りの非通常剰余は
定理\ref{thm:rank-fifteen-pencil-infinity}、命題\ref{prop:rank-fifteen-octic-residues}、
定理\ref{thm:rank-fifteen-octic-boundary}で認証済みである。完備曲線の各代数点は
109上の付値で還元でき、超越点は幾何学的一般点からの体拡大である。
通常な窓の一般Hilbert次元の下界に四中心商による上界を合わせ、
定理\ref{thm:rank-fourteen-all-ordinary}の自由性と左右Koszul複体の議論を適用する。
\end{proof}
